\section{文件管理}

\subsection{文件系统的基本概念}

\begin{mydef}{文件与文件系统}{file-system}
\textbf{文件}是具有符号名的一组相关信息的集合，是外存中数据组织的基本单位。\textbf{文件系统}是 OS 中负责文件存储、检索、保护和共享的子系统。
\end{mydef}

\begin{conclusion}{文件的逻辑结构}{file-logical-struct}
\begin{itemize}
    \item \textbf{无结构文件（流式文件）}：以字节流为单位（如 UNIX/DOS 文件），灵活性最高。
    \item \textbf{有结构文件（记录式文件）}：由若干记录组成。分为定长记录和变长记录。
\end{itemize}
\end{conclusion}

\subsection{文件的物理结构（文件的分配方式）}

\begin{conclusion}{三种文件物理结构}{file-physical}
\begin{tablebox}{文件分配方式对比}
\centering
\small
\begin{tabularx}{\linewidth}{|p{0.20\linewidth}|X|X|X|p{0.15\linewidth}|}
\hline
分配方式 & 实现 & 优点 & 缺点 & FAT \\
\hline
连续分配 & 每个文件占连续盘块 & 顺序访问快 & 有外部碎片，难扩展 & — \\
链接分配（隐式）& 每块含下一块指针 & 无外部碎片 & 随机访问慢 & — \\
链接分配（显式/FAT）& FAT 集中存指针 & FAT 在内存时，遍历比隐式链接快；无外部碎片 & 定位第 $i$ 块仍需沿链逐项查找；FAT 占内存 & FAT12/16/32 \\
索引分配 & 索引块集中存放数据块号 & 支持直接访问；多级索引可支持大文件 & 索引块有空间开销 & UNIX i-node \\
\hline
\end{tabularx}
\end{tablebox}

\textbf{UNIX i-node 混合索引：}
\begin{itemize}
    \item 直接块（通常 10-12 个）：存文件数据块号，小文件直接访问。
    \item 一级间接块：存指向数据块的指针（指向一个块，该块中存数据块号）。
    \item 二级间接块：存指向一级间接块的指针。
    \item 三级间接块：存指向二级间接块的指针。
\end{itemize}
\textbf{408 常考：}给定块大小和指针大小，计算 i-node 支持的最大文件大小。
\end{conclusion}

\begin{center}
\begin{tikzpicture}[>=stealth, font=\scriptsize]
% ===== 连续分配 =====
\node[anchor=west, text=Burgundy] at (0,3.2) {\textbf{连续分配}};
\foreach \i/\x in {0/2.6,1/3.7,2/4.8,3/5.9} {
  \node[draw, rectangle, minimum width=0.95cm, minimum height=0.55cm, fill=SoftBlue] (c\i) at (\x,3.2) {块\i};
}
\draw[decorate, decoration={brace, amplitude=3pt, mirror}, thick]
  ([yshift=-0.35cm]c0.south) -- ([yshift=-0.35cm]c3.south);
\node[anchor=west, align=left] at (7.2,3.2) {\scriptsize 文件占用连续盘块\\[1pt]顺序访问快};
% ===== 链接分配 =====
\node[anchor=west, text=Burgundy] at (0,1.0) {\textbf{链接分配}};
\foreach \i/\x in {0/2.6,1/4.1,2/5.6,3/7.1} {
  \node[draw, rectangle, minimum width=1.2cm, minimum height=0.55cm, fill=SoftTeal] (l\i) at (\x,1.0) {块\i};
}
\draw[->] (l0) -- (l1);
\draw[->] (l1) -- (l2);
\draw[->] (l2) -- (l3);
\node[anchor=west, align=left] at (8.6,1.0) {\scriptsize 每块含指向下一\\[1pt]块的指针（隐式）};
% ===== 索引分配 =====
\node[anchor=west, text=Burgundy] at (0,-1.2) {\textbf{索引分配}};
\node[draw, rectangle, minimum width=1.4cm, minimum height=0.65cm, fill=SoftAmber] (idx) at (2.6,-1.2) {索引块};
\node[draw, rectangle, minimum width=0.95cm, minimum height=0.55cm, fill=SoftAmber!55] (d0) at (4.6,-0.45) {块0};
\node[draw, rectangle, minimum width=0.95cm, minimum height=0.55cm, fill=SoftAmber!55] (d1) at (4.6,-1.2) {块1};
\node[draw, rectangle, minimum width=0.95cm, minimum height=0.55cm, fill=SoftAmber!55] (d2) at (4.6,-1.95) {块2};
\draw[->] (idx.east) -- (d0.west);
\draw[->] (idx.east) -- (d1.west);
\draw[->] (idx.east) -- (d2.west);
\node[anchor=west, align=left] at (6.4,-1.2) {\scriptsize 索引块集中存数据\\[1pt]块号，支持直接访问};
\end{tikzpicture}
\end{center}

\begin{center}
\begin{tikzpicture}[>=stealth, font=\scriptsize]
\node[draw, rectangle, fill=SoftGray, minimum width=2.4cm, minimum height=0.6cm] (e1) at (-8.0,3.0) {直接块$\times$10};
\node[draw, rectangle, fill=SoftGray, minimum width=2.4cm, minimum height=0.6cm] (e2) at (-8.0,1.2) {一级间接};
\node[draw, rectangle, fill=SoftGray, minimum width=2.4cm, minimum height=0.6cm] (e3) at (-8.0,-0.6) {二级间接};
\node[draw, rectangle, fill=SoftGray, minimum width=2.4cm, minimum height=0.6cm] (e4) at (-8.0,-2.4) {三级间接};
\node[font=\scriptsize, text=Navy] at (-8.0,3.7) {\textbf{i-node（索引节点）}};
\node[draw, rectangle, fill=SoftBlue, minimum width=2.6cm, minimum height=0.6cm] (d1) at (3.5,3.0) {数据块};
\node[draw, rectangle, fill=SoftTeal, minimum width=2.6cm, minimum height=0.6cm] (s1) at (-2.4,1.2) {一级间接块};
\node[draw, rectangle, fill=SoftBlue, minimum width=2.6cm, minimum height=0.6cm] (d2) at (3.5,1.2) {数据块};
\node[draw, rectangle, fill=SoftTeal, minimum width=2.6cm, minimum height=0.6cm] (b2) at (-5.0,-0.6) {二级间接块};
\node[draw, rectangle, fill=SoftTeal, minimum width=2.6cm, minimum height=0.6cm] (s2) at (-0.8,-0.6) {一级间接块};
\node[draw, rectangle, fill=SoftBlue, minimum width=2.6cm, minimum height=0.6cm] (d3) at (3.5,-0.6) {数据块};
\node[draw, rectangle, fill=SoftTeal, minimum width=2.6cm, minimum height=0.6cm] (t3) at (-5.6,-2.4) {三级间接块};
\node[draw, rectangle, fill=SoftTeal, minimum width=2.6cm, minimum height=0.6cm] (b3) at (-2.8,-2.4) {二级间接块};
\node[draw, rectangle, fill=SoftTeal, minimum width=2.6cm, minimum height=0.6cm] (s3) at (0.0,-2.4) {一级间接块};
\node[draw, rectangle, fill=SoftBlue, minimum width=2.6cm, minimum height=0.6cm] (d4) at (3.5,-2.4) {数据块};
\draw[->] (e1.east) -- (d1.west);
\draw[->] (e2.east) -- (s1.west);
\draw[->] (s1.east) -- (d2.west);
\draw[->] (e3.east) -- (b2.west);
\draw[->] (b2.east) -- (s2.west);
\draw[->] (s2.east) -- (d3.west);
\draw[->] (e4.east) -- (t3.west);
\draw[->] (t3.east) -- (b3.west);
\draw[->] (b3.east) -- (s3.west);
\draw[->] (s3.east) -- (d4.west);
\node[font=\scriptsize, MutedText] at (-0.9,-3.5) {直接块数决定小文件无需间接层；文件越大，间接层越多，可支持的文件上限越大};
\end{tikzpicture}
\end{center}

\begin{attention}
\textbf{408 重点辨析：FCB（文件控制块）与 i-node。} FCB 包含文件名和索引指针等信息。UNIX 将文件名和索引指针分离：目录项只存文件名和 i-node 号，i-node 存除文件名外的所有信息。这样做的好处是支持文件硬链接（多个目录项指向同一 i-node）。
\end{attention}

\subsection{目录结构}

\begin{conclusion}{目录结构演化}{directory-struct}
\begin{itemize}
    \item \textbf{单级目录}：所有文件在一个目录下。不允许重名，查找慢。
    \item \textbf{两级目录}：主文件目录（MFD）+ 用户文件目录（UFD）。不同用户可有同名文件。
    \item \textbf{树形目录（多级目录）}：目录以树形组织。使用\textbf{绝对路径}（/a/b/file）和\textbf{相对路径}（当前目录为基准）。
    \item \textbf{无环图目录}：允许一个文件或目录出现在多个父目录下（共享）。需引入引用计数防止删除正在使用的文件。
\end{itemize}
\end{conclusion}

\begin{center}
\begin{tikzpicture}[>=stealth, font=\scriptsize]
\node[draw, rectangle, fill=SoftGray, minimum width=1.8cm, minimum height=0.55cm] (root) at (0,0) {/};
\node[draw, rectangle, fill=SoftBlue, minimum width=2.0cm, minimum height=0.5cm] (prog) at (-2.4,-1.1) {Program};
\node[draw, rectangle, fill=SoftBlue, minimum width=2.0cm, minimum height=0.5cm] (doc) at (0.8,-1.1) {Document};
\node[draw, rectangle, fill=SoftTeal, minimum width=2.0cm, minimum height=0.5cm] (java) at (-3.2,-2.3) {Java-Prog};
\node[draw, rectangle, fill=SoftAmber, minimum width=2.0cm, minimum height=0.5cm] (f1) at (-1.2,-2.3) {f1.java};
\node[draw, rectangle, fill=SoftAmber, minimum width=2.0cm, minimum height=0.5cm] (fdoc) at (0.8,-2.3) {报告.pdf};
\draw[->] (root) -- (prog);
\draw[->] (root) -- (doc);
\draw[->] (prog) -- (java);
\draw[->] (prog) -- (f1);
\draw[->] (doc) -- (fdoc);
\node[font=\scriptsize, black!60] at (-0.6,-3.25) {绝对路径 /Program/f1.java；相对路径（当前在 /Program）：f1.java};
\end{tikzpicture}
\end{center}

\subsection{文件存储空间管理}

\begin{conclusion}{外存空闲空间管理方法}{free-space}
\begin{itemize}
    \item \textbf{空闲表法}：类似内存动态分区，维护空闲盘块表。分配用首次适应等算法。
    \item \textbf{空闲链表法}：所有空闲盘块连成链表。分配/回收效率低（需 I/O）。
    \item \textbf{位示图（Bitmap）}：每位对应一个盘块（0=空闲，1=已分配）。$m$ 个盘块需 $m$ bit。\textbf{408 常考：}给定盘块号，计算在位示图中的行列号。
    \item \textbf{成组链接法（UNIX 采用）}：空闲盘块分组管理，每组第一块存下一组空闲块的块号和块数。兼顾效率和空间。
\end{itemize}
\end{conclusion}

\subsection{文件共享与保护}

\begin{conclusion}{文件共享}{file-sharing}
\begin{itemize}
    \item \textbf{硬链接（Hard Link）}：多个目录项指向同一 i-node。无法跨文件系统，不能链接目录（防止循环）。
    \item \textbf{软链接（Symbolic Link）}：独立文件，存放目标文件路径。可跨文件系统、可链接目录。删除原文件后软链接失效（悬空链接）。
\end{itemize}
\end{conclusion}

\begin{center}
\begin{tikzpicture}[>=stealth, font=\scriptsize]
\node[draw, rectangle, fill=SoftGray, minimum width=1.6cm, minimum height=0.7cm] (ino) at (0,0) {i-node};
\node[draw, rectangle, fill=SoftBlue, minimum width=1.9cm, minimum height=0.5cm] (dirA) at (-2.0,1.3) {目录项A};
\node[draw, rectangle, fill=SoftBlue, minimum width=1.9cm, minimum height=0.5cm] (dirB) at (2.0,1.3) {目录项B};
\draw[->] (dirA) -- (ino);
\draw[->] (dirB) -- (ino);
\node[font=\scriptsize, black!70] at (0,2.1) {硬链接：多个目录项指向同一 i-node};
\node[font=\scriptsize, black!60] at (0,0.78) {（引用计数；不能跨文件系统）};
\node[draw, rectangle, fill=SoftAmber, minimum width=2.2cm, minimum height=0.5cm] (link) at (-0.6,-1.5) {软链接（独立文件）};
\node[draw, rectangle, fill=SoftGray, minimum width=2.2cm, minimum height=0.5cm] (data) at (3.6,-1.5) {目标文件};
\draw[->] (link.east) -- node[above,font=\scriptsize]{存放路径} (data.west);
\node[font=\scriptsize, black!70] at (1.4,-2.5) {软链接：可跨文件系统，目标删除后悬空};
\end{tikzpicture}
\end{center}

\subsection{磁盘调度算法}

\begin{conclusion}{磁盘访问时间与调度}{disk-scheduling}
磁盘访问时间 $T = T_s$（寻道时间）+ $T_r$（旋转延迟）+ $T_t$（传输时间）。

\textbf{寻道时间远大于后两者}，因此调度算法以优化寻道为目标。

\begin{tablebox}{磁盘调度算法对比}
\centering
\begin{tabularx}{\linewidth}{|l|X|X|X|}
\hline
算法 & 策略 & 优点 & 缺点 \\
\hline
FCFS & 按请求顺序 & 公平 & 性能差 \\
SSTF & 最短寻道优先 & 性能好 & 饥饿 \\
SCAN（电梯算法）& 磁头单向移动，到端点折返 & 无饥饿 & 端点附近不公平 \\
C-SCAN & 单向服务，到端点跳回起点 & 等待时间均匀 & — \\
LOOK/C-LOOK & 类似 SCAN/C-SCAN，但到最远请求即折返 & 更实用 & — \\
\hline
\end{tabularx}
\end{tablebox}

\textbf{408 常考：}给定磁盘请求序列和当前磁头位置/方向，用 SSTF/SCAN/C-SCAN 画出磁头移动轨迹，计算总寻道长度和平均寻道长度。
\end{conclusion}

\subsection{408 高频考点速查}

\begin{conclusion}{文件管理 — 408常考点}{os4-keypoints}
\begin{enumerate}
    \item \textbf{i-node 混合索引}：最大文件长度计算。
    \item \textbf{位示图}：盘块号 $\leftrightarrow$ 字号/位号的换算（$b=n\times\text{字长}+m$）。
    \item \textbf{磁盘调度}：SSTF/SCAN/C-SCAN 的磁头移动序列和总寻道长度。
    \item \textbf{硬链接 vs 软链接}：引用计数、跨文件系统、删除行为。
    \item \textbf{FAT 文件系统}：FAT 表大小计算（给定盘块数和 FAT 表项位数）。
\end{enumerate}
\end{conclusion}

\newpage
\section{习题}

\subsection{基础过关题}

\begin{enumerate}

\item \textbf{（真题风格·选择题）} 某文件系统的根目录 \texttt{/} 下有子目录 \texttt{Program} 和 \texttt{Document}；\texttt{Program} 下有子目录 \texttt{Java-Prog} 以及文件 \texttt{f1.java}。当前工作目录为 \texttt{/Program}，则文件 \texttt{f1.java} 的绝对路径名为
\begin{center}
\begin{tabular}{ll}
(A) \texttt{/Program/f1.java} & (B) \texttt{/Document/Program/f1.java} \\[6pt]
(C) \texttt{Program/f1.java} & (D) \texttt{/Java-Prog/f1.java}
\end{tabular}
\end{center}

\ifshowsolutions\par\noindent\textbf{解析：} f1.java在Program目录下，全路径\texttt{/Program/f1.java}。\boxed{\textbf{答案：}(A)}\else \answerarea \fi

\vspace{8pt}

\item \textbf{（真题风格·选择题）} 在磁盘调度算法中，通常称为"电梯算法"的是
\begin{center}
\begin{tabular}{ll}
(A) FCFS & (B) SSTF \\[6pt]
(C) SCAN & (D) C-SCAN
\end{tabular}
\end{center}

\ifshowsolutions\par\noindent\textbf{解析：} SCAN算法磁头像电梯一样单向移动直到端点折返，故称电梯算法。答案(C)。\boxed{\textbf{答案：}(C)}\else \answerarea \fi

\vspace{8pt}

\item \textbf{（真题风格·选择题）} 假设磁盘有300个磁道（0-299），磁头当前在120号磁道并向磁道号减小的方向移动。请求序列为：90, 70, 150, 60, 80, 20。用SCAN算法（LOOK变体），磁头移动的总磁道数为
\begin{center}
\begin{tabular}{ll}
(A) 140 & (B) 200 \\[6pt]
(C) 230 & (D) 280
\end{tabular}
\end{center}

\ifshowsolutions\par\noindent\textbf{解析：} LOOK 向减小方向服务至最小请求 20 后折返：$120\to90\to80\to70\to60\to20\to150$。总移动距离为 $30+10+10+10+40+130=230$ 个磁道。\boxed{\textbf{答案：}(C)}\else \answerarea \fi

\vspace{8pt}

\item \textbf{（真题风格·选择题）} 在文件系统中，文件控制块（FCB）的作用是
\begin{center}
\begin{tabular}{ll}
(A) 存放文件的数据内容 & (B) 存放文件的属性信息 \\[6pt]
(C) 存放文件的目录结构 & (D) 存放文件的访问权限列表
\end{tabular}
\end{center}

\ifshowsolutions\par\noindent\textbf{解析：} FCB存放文件的属性信息（文件名、大小、创建时间、物理位置等）。目录由FCB组成。UNIX将文件名从FCB分离到目录项，FCB变为i-node。\boxed{\textbf{答案：}(B)}\else \answerarea \fi

\vspace{8pt}

\item \textbf{（基础题）} 下列选项中，属于文件系统的基本功能的是
\begin{center}
\begin{tabular}{ll}
(A) 文件按名存取 & (B) 进程调度 \\[6pt]
(C) 处理机分配 & (D) 虚拟内存地址变换
\end{tabular}
\end{center}

\ifshowsolutions\par\noindent\textbf{解析：}文件目录把文件名映射到 FCB/i-node，实现按名存取；其余选项属于处理机或内存管理。\boxed{\textbf{答案：}(A)}\else \answerarea \fi

\vspace{8pt}

\item \textbf{（计算题）} 某磁盘有8000个柱面（0--7999），磁头当前在150号柱面并向柱面号增大的方向移动。请求序列为：86, 147, 913, 1774, 948, 1509, 1022, 1750, 130。分别用 SSTF 和 LOOK 算法，计算磁头移动的总柱面数。

\ifshowsolutions\par\noindent\textbf{解析：}
SSTF：从 150 开始，访问顺序为 $150\to147\to130\to86\to913\to948\to1022\to1509\to1750\to1774$，各段移动量依次为 3、17、44、827、35、74、487、241、24，总移动量为 1752 个柱面。

LOOK：先向增大方向依次访问 $913,948,1022,1509,1750,1774$，从 150 到 1774 共移动 1624 个柱面；再折返依次访问 $147,130,86$，从 1774 到 86 共移动 1688 个柱面。因此总移动量为 $1624+1688=3312$ 个柱面。\else \answerarea \fi

\end{enumerate}

\subsection{真题风格与综合训练}

\begin{enumerate}[resume]

\item \textbf{（真题风格·综合题）} 某UNIX文件系统i-node包含8个直接块指针、1个一级间接块指针、1个二级间接块指针。盘块大小为4 KB，每个指针占4字节。
\begin{enumerate}
    \item[(1)] 该文件系统支持的最大文件大小是多少？
    \item[(2)] 访问文件偏移量 5 MiB 处的数据，需要访问几次磁盘？
    \item[(3)] 若文件系统改用FAT32管理磁盘空间，FAT表项32位，磁盘总容量为4 GB，盘块大小4 KB。FAT表需占用多少存储空间？
\end{enumerate}

\ifshowsolutions\par\noindent\textbf{解析：}
(1) 直接块：$8\times 4\text{ KiB}=32\text{ KiB}$。一级间接：每块可存$4\text{ KiB}/4\text{ B}=1024$个指针，共$1024\times 4\text{ KiB}=4\text{ MiB}$。二级间接：$1024^2\times 4\text{ KiB}=4\text{ GiB}$。最大文件大小为 $4\text{ GiB}+4\text{ MiB}+32\text{ KiB}=4{,}299{,}194{,}368$ B，约 $4.00394\text{ GiB}$。

(2) 5 MiB 已超过直接块和一级间接区合计的 $32\text{ KiB}+4\text{ MiB}$，位于二级间接区。若相关块均未缓存，需要读取二级间接根块、其指向的下一级索引块和目标数据块，共 3 次磁盘 I/O。

(3) 盘块数$=4\text{ GB}/4\text{ KB}=1\text{M}$块。FAT表项32位=4 B，FAT大小$=1\text{M}\times 4\text{ B}=4\text{ MB}$。\else \answerarea \fi

\vspace{8pt}

\item \textbf{（真题风格·综合题）} 某系统采用位示图法管理磁盘空闲空间。磁盘共有$2^{24}$个盘块，字长为32位。
\begin{enumerate}
    \item[(1)] 位示图需要多少个字？
    \item[(2)] 给出盘块号$b$到位示图中字号$i$和位号$j$的换算公式；
    \item[(3)] 若盘块号500000在位示图中对应字号为15625，验证(2)中公式的正确性。
\end{enumerate}

\ifshowsolutions\par\noindent\textbf{解析：}
(1) $2^{24}\div 32=2^{24}\div 2^{5}=2^{19}=524288$个字。

(2) $i=\lfloor b/32\rfloor,\ j=b\bmod 32$。

(3) $500000\div 32=15625$（字号$i=15625$），$500000\bmod 32=0$（位号$j=0$）。验证：$15625\times 32+0=500000$。\else \answerarea \fi

\end{enumerate}

\vspace{8pt}
\noindent\textbf{拓展阅读：}
\begin{itemize}
    \item OSTEP 第39-40章——文件系统实现与i-node详解。
    \item 陈海波《现代操作系统：原理与实现》第7章——Ext4与F2FS文件系统。
\end{itemize}
